Federated learning over consumer IoT devices is increasingly subject to
regulatory regimes (GDPR Article~17, EU AI Act Article~15, CCPA
\S~1798.105) that require data subjects to withdraw consent and have
the trained model demonstrably forget their contribution. Existing work
treats consent management, machine unlearning, and on-chain auditability
as three disjoint research threads. We close the
\emph{consent--unlearning--verification triangle} with \sysname, a
protocol that binds the three layers in a single chain-enforced
pipeline. Our contributions are: (i) a Consent Lifecycle State Machine
(\clsm) modelling consent as a 6-state finite automaton implemented as
soulbound ERC-721 tokens; (ii) a Verifiable Unlearning Protocol (\vup)
that gates on-chain acceptance on three independent membership
inference attacks; (iii) a three-contract Solidity architecture
deployed on PureChain (gas-free, PoA$^2$) and Polygon Amoy (gas-based,
PoS) for the first apples-to-apples cost comparison of consent-aware
federated learning under both economic models; and (iv) an end-to-end
dual-domain empirical evaluation on a 9-client N-BaIoT autoencoder
federation (primary) cross-validated on a 5-client CASAS BiLSTM
activity classifier, across three approximate-unlearning paradigms
plus an exact-retraining baseline. We
observe per-data-subject lifecycle costs of $\$0$ on PureChain versus
\$0.0012 on Polygon Amoy and per-round on-chain validation costs that
scale linearly with the participant count $N$, reaching \$0.0176 at
$N=500$. The protocol accepts approximate-unlearning attestations
($\text{MIA}\!<\!0.55$) while rejecting both naive
retraining (whose distributional uniqueness raises false-positive MIA
scores) and any attestation lacking a 3-of-5 verifier threshold,
demonstrating that the on-chain gate is faithful to MIA evidence rather
than to the data controller's claim.
